\section{Flat-Field Generation and Analysis Configuration} \label{app:config}

For reproducibility we record here the exact inputs used to build the deep $i$-band combined
flat and to run the single-detector analysis.

The analysis notebooks read the combined flats directly from the Butler at the USDF with the
following configuration:
%
{\footnotesize
\begin{verbatim}
BUTLER_PATH = "/repo/main"
COLLECTIONS = "u/plazas/LSSTCam/calib/none/treerings/flatGen-i.20260319b"
BAND        = "i"
\end{verbatim}
}
%
The chromaticity analysis instead reads the five-band combined flats from the standard
\texttt{LSSTCam/calib} collection.

The combined $i$-band flat was produced from 290 individual PTC flat exposures (taken with the
3s\_v2 ``natural-frequency'' sequencer on 2025-12-28) using the LSST Science Pipelines
flat-generation pipeline, submitted through the Batch Processing Service (BPS) at the USDF.
The submission script is reproduced below; the \texttt{SELECTION\_FLAT\_i} exposure list is
abbreviated here (the full set of 290 exposure IDs is preserved with the analysis code).
%
{\footnotesize
\begin{verbatim}
# Generate a combined flat image from i-band PTC flats to study tree rings (2026-02-26)
# Run this at USDF.

# STEP 0: set up the LSST Science Pipelines stack
source /sdf/group/rubin/sw/w_latest/loadLSST.bash
setup lsst_sitcom

# STEP 1: work from the home directory
cd ${HOME}

# STEP 2: environment / data selection
export USER_CALIB_PREFIX="u/plazas/"
export INSTRUMENT=LSSTCam
export TICKET="none"
export REPO=main
export RAW_COLLECTION=LSSTCam/raw/all
# Curated calibration inputs:
#   LSSTCam/calib/DM-53722: defects, bias, dark, linearity, etc.
#   LSSTCam/calib/DM-52726: flat gradient reference
#   LSSTCam/calib/DM-51599: CBP crosstalk for 3s_v1
#   LSSTCam/calib/DM-51669: curated calibrations
#   LSSTCam/calib/photodiode: photodiode data
export CALIB_COLLECTIONS=LSSTCam/calib/DM-53722,LSSTCam/calib/DM-52726,\
LSSTCam/calib/DM-51599,LSSTCam/calib/DM-51669,LSSTCam/calib/photodiode
export TAG=treerings
export RERUN=20260319b   # bump to "c", "d", ... to rerun

# Flats: 290 i-band PTC flats, 3s_v2 ("natural frequency") sequencer.
# Detector 122 is excluded.
export SELECTION_FLAT_i="instrument='LSSTCam' and detector not in (122) and \
exposure in (2025122800010, 2025122800011, 2025122800012, 2025122800013,
             2025122800014, 2025122800015, 2025122800024, 2025122800025,
             ...   # 290 exposure IDs in total (full list with the analysis code)
             2025122801080, 2025122801081, 2025122801084, 2025122801085)"

# STEP 3: sanity-check the resolved BPS YAML
cat $CP_PIPE_DIR/bps/templates/bps_flat_i.yaml | envsubst

# STEP 4: submit
bps submit -i ${RAW_COLLECTION},${CALIB_COLLECTIONS} \
    $CP_PIPE_DIR/bps/templates/bps_flat_i.yaml
\end{verbatim}
}
%
The per-exposure Poisson noise averages down as $\sqrt{290}$ in the combine, which is what
brings the $\sim10^{-4}$ tree-ring modulation above the single-flat noise floor.

\section{Per-Detector Tree-Ring Profile Tables} \label{app:tables}

We provide tabulated radial tree-ring profiles for all detectors and bands. These tables are available online at [DOI TBD] and include:

\begin{itemize}
\item Detector ID
\item Band
\item Fitted center $(x_0, y_0)$
\item Initial center from HyeYun's catalog $(x_{0,\rm init}, y_{0,\rm init})$
\item Radial coordinate $r$ (pixels)
\item Profile amplitude $w(r)$ (fractional units)
\item Uncertainty $\sigma_w(r)$
\item Mask fraction versus $r$
\item Quality flags
\end{itemize}

A sample of the table format is shown below.

\begin{deluxetable*}{cccccccc}
\tablecaption{Sample Tree-Ring Profile Table\label{tab:sample_profile}}
\tablewidth{0pt}
\tablehead{
\colhead{Detector} & \colhead{Band} & \colhead{$r$} & \colhead{$w(r)$} & \colhead{$\sigma_w$} & \colhead{Mask Frac.} & \colhead{Counts} & \colhead{Flag} \\
\colhead{} & \colhead{} & \colhead{(pix)} & \colhead{} & \colhead{} & \colhead{} & \colhead{} & \colhead{}
}
\startdata
R02\_S02 & i & 100.5 & 0.0012 & 0.0003 & 0.02 & 1245 & 0 \\
R02\_S02 & i & 101.5 & 0.0015 & 0.0003 & 0.02 & 1250 & 0 \\
\dots & \dots & \dots & \dots & \dots & \dots & \dots & \dots
\enddata
\tablecomments{Full tables for all detectors and bands are available online.}
\end{deluxetable*}

%%%%%%%%%%%%%%%%%%%%%%%%%%%%%%%%%%%%%%%%%%%%%%%%%%%%%%%%%%%%%%%%%%%%%%%%%%%%%%%%
